\begin{table}[htbp] \centering \caption{Summary Statistics China, Balanced Panel}\label{tab:summary_stats_CHN_bal} \begin{threeparttable} \begin{tabular}{l cccc} \toprule
& All & Rural & Urban & Difference \\
 &  &  &  & $t$-test \\
\midrule
Location &  &   56.2\% &   43.8\% &  \\  \addlinespace \addlinespace
Log Consumption & 10.39 & 10.21 & 10.63 & -0.43*** \\
 & (0.92) & (0.89) & (0.89) &  \\  \addlinespace
Log Income & 8.60 & 8.08 & 9.13 & -1.05*** \\
 & (1.91) & (2.06) & (1.58) &  \\  \addlinespace
Female & 0.52 & 0.51 & 0.53 & -0.01*** \\
 & (0.50) & (0.50) & (0.50) &  \\  \addlinespace
Age (years) & 49.28 & 49.34 & 49.20 & 0.13 \\
 & (14.61) & (14.38) & (14.91) &  \\  \addlinespace
Education (years) & 6.92 & 5.70 & 8.49 & -2.79*** \\
 & (4.78) & (4.46) & (4.73) &  \\  \addlinespace
Household Size & 4.12 & 4.39 & 3.77 & 0.62*** \\
 & (1.81) & (1.87) & (1.67) &  \\  \addlinespace


\cmidrule{2-5}
Observations & 56,855 & 31,968 & 24,887 &  \\  
Individuals & 14,214 &  &  &  \\
Non-switchers &   91.8\% &  &  &  \\  \addlinespace
\bottomrule \end{tabular} \begin{tablenotes}[flushleft] \footnotesize \item{Summary statistics for China for the balanced panel across all waves. Source: China survey. The table reports means and standard deviations (in parentheses) based on individual-year pairs. See section 3 for further details. All variables have the same number of observations, except for income, which is missing for some observations. Income has 25,530 observations.} \end{tablenotes} \end{threeparttable} \end{table}